\subsection{トピックと参考文献の対応}
この表は、SWEBOK が示す「トピック対参考資料の行列」を、学習用に簡略化したものである。詳しく学ぶときは、各トピックに対応する文献を読む。

\par\smallskip\noindent\refstepcounter{table}\label{tab:ch02-07}表 \thetable: トピックと参考文献の対応におけるトピック・主な参考資料\par\smallskip
表 \thetable は、トピックと参考文献の対応におけるトピック・主な参考資料を比較・確認しやすい形で整理したものである。\par\smallskip
\begin{longtable}{P{0.42\linewidth}P{0.50\linewidth}}
\toprule
トピック & 主な参考資料 \\
\midrule
1. ソフトウェアアーキテクチャの基礎 & Bass et al.; Maier and Rechtin; Rozanski and Woods; Taylor et al. \\
1.1 「アーキテクチャ」という語の意味 & Bass et al.; Budgen; Maier and Rechtin \\
1.2 利害関係者と関心事 & Bass et al.; Rozanski and Woods; Taylor et al. \\
1.3 アーキテクチャの用途 & Bass et al.; Rozanski and Woods \\
2. アーキテクチャ記述 & Bass et al.; Rozanski and Woods; Sommerville; Taylor et al. \\
2.1 ビューとビューポイント & Budgen; Maier and Rechtin; Rozanski and Woods; Sommerville \\
2.2 パターン、様式、参照アーキテクチャ & Bass et al.; Budgen; Rozanski and Woods; Sommerville; Taylor et al. \\
2.3 アーキテクチャ記述言語と枠組み & Bass et al.; Maier and Rechtin; Rozanski and Woods; Taylor et al. \\
2.4 重要な意思決定としてのアーキテクチャ & Rozanski and Woods; Sommerville \\
3. アーキテクチャ過程 & Maier and Rechtin; Rozanski and Woods; Taylor et al. \\
3.1 文脈の中のアーキテクチャ & Maier and Rechtin; Taylor et al. \\
3.1.1 アーキテクチャと設計の関係 & Sommerville; Taylor et al. \\
3.2 アーキテクチャ設計 & Bass et al. \\
3.2.1 アーキテクチャ分析 & Bass et al.; Taylor et al. \\
3.2.2 アーキテクチャ統合 & Bass et al. \\
3.2.3 アーキテクチャ評価 & Bass et al.; Rozanski and Woods \\
3.3 実践、方法、戦術 & Bass et al.; Maier and Rechtin; Rozanski and Woods \\
3.4 大規模なアーキテクチャ活動 & Maier and Rechtin; Sommerville \\
4. アーキテクチャ評価 & Bass et al.; Rozanski and Woods; Taylor et al. \\
4.1 アーキテクチャの「よさ」 & Bass et al.; Budgen \\
4.2 アーキテクチャについての推論 & Rozanski and Woods \\
4.3 アーキテクチャレビュー & Bass et al. \\
4.4 アーキテクチャ測度 & Bass et al. \\
\bottomrule
\end{longtable}
